\section{Conclusions}
\label{sec:conclusions}

\todo{Write conclusions based on final results. The following structure is suggested:}

This project analyzed the environmental impact of disposable electronic cigarettes and applied Design for Disassembly principles to redesign the product for improved end-of-life management. The key findings are:

\begin{enumerate}
    \item \todo{Finding 1: Current disposable e-cigarettes are designed as sealed, single-use units that prevent any material recovery. The LCA showed that [greatest impact category] is dominated by [SBOM input / life cycle stage].}

    \item \todo{Finding 2: Applying DfD principles---particularly snap-fit assembly, modular architecture, and prioritized battery accessibility---reduced disassembly time from X minutes to under 30 seconds and eliminated the need for destructive disassembly.}

    \item \todo{Finding 3: The proposed Universal Removable Battery Standard, defining three standardized battery sizes with mechanical and electrical interface specifications, could reduce per-device environmental impacts by X\% across [impact categories] by enabling battery reuse across multiple device lifecycles.}

    \item \todo{Finding 4: The alternative concept with mono-material PP housing and modular design achieves material recovery rates exceeding 90\% for critical metals, compared to approximately 0\% for the current design.}
\end{enumerate}

The proposed standard aligns with the EU Batteries Regulation (2023/1542), which will require removable and replaceable batteries in all portable electronics by February 2027. Early adoption of such a standard would provide manufacturers with a compliance pathway while addressing the urgent environmental and safety concerns associated with the current generation of disposable e-cigarettes.

\subsection{Limitations}

\todo{Discuss limitations:}
\begin{itemize}
    \item LCA assumptions and data quality
    \item Idealized recycling rate assumptions
    \item Consumer behavior uncertainty (will users actually remove batteries?)
    \item Standardization governance (who administers the standard?)
    \item Potential for the standard to extend the lifecycle of nicotine products (public health consideration)
\end{itemize}

\subsection{Future Work}

\todo{Suggest future directions:}
\begin{itemize}
    \item Prototype fabrication and physical disassembly testing
    \item Consumer behavior studies on battery removal compliance
    \item Detailed cost--benefit analysis for manufacturers
    \item Engagement with standards bodies (IEC, ISO) for formal standardization
    \item Extension to rechargeable/refillable e-cigarette designs
\end{itemize}
